\scanpage{132}
\mspage{65}
contre il y en a entre $\Pi_{d,d'}$ et $\Pi_d$.
Précisons ce point, en considérant un
morphisme
\[
(215)\qquad\Delta_r\xrightarrow{\alpha}\Delta_{r'}\qquad\text{\uuline{tel que}}\quad\alpha(0)=0
\]
alors $D_{\Delta_{r'}}\longrightarrow D_{\Delta_r}$ est un morphisme
étale fini, donc aussi $D^*_{\Delta_{r'}}\unsure{\text{---}|\longrightarrow}D^*_{\Delta_r}$ induit
au dessus de $D^*_{\Delta_r}$) donc on trouve
un diagramme $2$-cartésien (qui \deleted{\unsure{suit|vient}} \inserted{peut}
\deleted{(209)} aussi se déduire de (209)
\diagram{\begin{tikzcd}[ampersand replacement=\&,column sep=large,row sep=large]
\&\widetilde\Pi_{\Delta_{r'}}\arrow[r]\arrow[d]\&\widetilde\Pi_{\Delta_r}\arrow[d]\\
(216)\&\Pi_{\Delta_{r'}}\arrow[r]\&\Pi_{\Delta_r}\arrow[ul,phantom,"2\text{-cart.}" description]
\end{tikzcd}}
)
\margin{Ceci montre que
les $\widetilde\Pi_{\Delta_r}\longleftarrow\Pi_{\Delta_r}$
et $\widetilde\Pi\unsure{\text{---}|\longrightarrow}\Pi$ sont
connus quand on
connaît les $\widetilde\Pi_d\longrightarrow\Pi_d$,
\unsure{des projections locales},
\deleted{\illegible{2}}
\uline{et} les \deleted{$\unsure{\Pi_d|P_d}$} \unsure{comme}
sur \deleted{$\unsure{\Pi_d|P_d}\longrightarrow\unsure{G_d|\Pi_d}$}
$D_d^{*!}$ sur $D_d$
\deleted{$\unsure{D_d}\longrightarrow\unsure{D^*_{\Delta_r}}$}}
où les flèches horizontales induisent sur
les $\pi_1$ des applications injectives à
image d'indice fini (on peut dire que
$\Pi_{\Delta_{r'}}\unsure{\text{---}|\longrightarrow}\Pi_{\Delta_r}$ fait de $\Pi_{\Delta_{r'}}$ un
rev. fini de $\Pi_{\Delta_r}$).

Toutes ces relations sont encore de
nature essentiellement formelle et triviale,
et marcheraient dans un site quelconque,
sans l'hypothèse spécifique des croisements
normaux, et sans que le remplacement des
\margin{131}
